\documentclass[14pt,a4paper]{extarticle}
\usepackage[margin=5em]{geometry}
\usepackage{amsthm,amssymb,amsfonts}
\usepackage{tikz}
\usetikzlibrary{shapes.misc}
\usepackage[skins]{tcolorbox}
\pagestyle{empty}
\newtcolorbox{Ex}[2][]{
                lower separated=false,
                colback=white!80!gray,
colframe=white, fonttitle=\bfseries,
colbacktitle=white!50!gray,
coltitle=black,
enhanced,
sharp corners,
#1}


\newtcolorbox{Thm}[2][]{
                lower separated=false,
                colback=white,
colframe=black,fonttitle=\bfseries,
colbacktitle=black,
coltitle=black,
enhanced,
sharp corners,
attach boxed title to top left={yshift=-0.1cm,xshift=-1em},
                 boxed title style={sharp corners,boxrule=0pt,colframe=white,colback=white,colbacktitle=black},
title=#2,#1}

\newtcolorbox{Exo}[2][]{
    shadow = {1.5mm}{-1.5mm}{0mm}{black!23!white},
    lower separated=false,
    colback=white,
colframe=black,fonttitle=\bfseries,
colbacktitle=black,
coltitle=black,
enhanced,
attach boxed title to top left={yshift=-0.2cm,xshift=1em},
boxed title style={boxrule=1pt,colframe=black,colback=white,colbacktitle=black
     },
,#1}


\begin{document}
\ 
\newpage
\begin{Thm}{Définition}
    Le dîner des talents ( D. D. T. ) est une soirée conviviale, rompant la monotonie de la pougne des préparationnaires au cours de laquelle les participants sont invités à mettre en scène un talent insoupçonné.
\end{Thm}


\begin{Ex}
    Exemple : faire du Rubik's cube avec ses pieds, rater un tour de magie\dots
\end{Ex}
\begin{Exo}
    \ \textbf{Point méthode }\\
    Pour assister au D. D. T. il suffit de se rendre au lycée Sainte Geneviève, 2 rue de l'école des postes, 78000 Versailles, au premier étage du bâtiment Saint-Joseph le mercredi 11 décembre 2024.
\end{Exo}
\smallskip

\textbf{Remarque fondamentale (donc HP) : } Une tenue chic est exigée (blouse blanche tolérée)


\begin{Thm}{Proposition}
    Les professeurs de la MP*1 sont invités au D. D. T. 
\end{Thm}

\textbf{Principe de démonstration. } Trivial. [LEE]

\begin{Thm}{Corollaire}
    M. Moulin est invité au D. D. T.
\end{Thm}
\textbf{Principe de démonstration. } Évident\\

\textbf{Remarque }En vue de la préparation de cet événement, il semble judicieux de priver (à grand regret) les MP*1 de leur devoir de Mathématiques hebdomadaire.\\

\begin{Exo}
\ \textbf{Exercice 1 } Présentez votre talent : 
\begin{enumerate}
    \item Identifiez un talent que en votre possession, si possible en rapport avec le thème.
    \item Préparez une démonstration pratique de ce talent, cette dernière doit être claire, rigoureuse, et peut comporter des éléments humoristique.
    \item Concluez.
\end{enumerate}

\end{Exo}

\end{document}